AIsafePy has a small, dependency-light core (\texttt{aisafepy.core})
plus three focused modules (\texttt{flow}, \texttt{stream},
\texttt{adapt}) and a \texttt{contrib} namespace that wraps upstream
classifiers. Every component produces or consumes the same
primitive, \texttt{GuardDecision}, so the modules compose end to end.
Figure~\ref{fig:architecture} shows the layering.

\begin{figure}[t]
  \centering
  \includegraphics[width=\linewidth]{figures/Figure1}
  \caption{AIsafePy architecture. The \texttt{core} layer defines
  the shared primitives. \texttt{flow} is the architectural defence
  via capability-based IFC at every tool-call boundary. \texttt{stream}
  is the detective layer with a tiered cascade. \texttt{adapt} closes
  the loop by compiling red-team failures into runtime guards.
  \texttt{contrib} adapts external classifiers to the Guard
  protocol.}
  \label{fig:architecture}
\end{figure}

\subsection{Core primitives}

\paragraph{GuardDecision.} Every guard returns a
\texttt{GuardDecision}, a frozen Pydantic model with the fields
\texttt{action} (allow / block / transform / escalate),
\texttt{confidence}, \texttt{tier}, \texttt{rationale},
\texttt{evidence}, \texttt{latency\_ms}, \texttt{severity}, and an
optional \texttt{fallback} or \texttt{transformed\_content}. This
structure addresses the observability problem in classifier-only
guardrails: developers can no longer be told that content was
unsafe without learning why, with what confidence, and at what
tier.

\paragraph{IFCViolation.} A specialized record used by
\texttt{aisafepy.flow}. Beyond the \texttt{GuardDecision} fields, it
carries the provenance set, capability set, integrity label, and a
reconstructed taint chain. An operator can answer the question
``which tool output ended up tainting this argument?'' from the
chain alone.

\paragraph{Budgets and progress.} The core ships iteration, token,
wall-clock, and dollar budgets. The budgets are thread-safe and emit
structured checkpoints at 25\%, 50\%, 75\%, and 100\% of capacity. A
\texttt{ProgressTracker} flags duplicate tool calls, cyclic
patterns, observation stagnation, and high semantic-similarity
loops, addressing the agent-loop pathology documented in public
bug trackers \citep{browseruse}.

\paragraph{Telemetry.} Every decision becomes an OpenTelemetry span
with the canonical attributes. The
library imports only \texttt{opentelemetry-api}. If the SDK is not
configured at runtime, a no-op tracer is used and attribute writes
are free.

\subsection{aisafepy.flow: capability-based IFC}

\paragraph{Lattice.} Each value crossing the IFC boundary is wrapped
in \texttt{Tainted[T]}, a generic record carrying three labels:

\begin{itemize}[leftmargin=2em, itemsep=0.1em]
  \item \texttt{provenance}: a frozen set of source identifiers,
    joined by union on composition.
  \item \texttt{capabilities}: a frozen set of capability tags,
    joined by union. The argument capability set must be a subset
    of the policy's declared ceiling for the tool.
  \item \texttt{integrity}: a label drawn from the three-element
    lattice TRUSTED $\succ$ UNTRUSTED $\succ$ QUARANTINED, joined
    by meet (worst wins).
\end{itemize}

For two Tainted values $a$ and $b$ produced by a binary composition:
\[
\mathrm{prov}(a \circ b) = \mathrm{prov}(a) \cup \mathrm{prov}(b),
\qquad
\mathrm{caps}(a \circ b) = \mathrm{caps}(a) \cup \mathrm{caps}(b),
\qquad
\mathrm{integ}(a \circ b) = \mathrm{integ}(a) \wedge \mathrm{integ}(b),
\]
where $\wedge$ is the meet of the integrity lattice.

\paragraph{Policy.} A \texttt{Policy} is a fluent builder that
declares three things: which sources produce which labels
(\texttt{label\_source}), which tools require which capabilities
and integrity floors (\texttt{require}), and which custom predicates
should outright deny a call (\texttt{deny\_if}). Listing~\ref{lst:policy}
shows the canonical setup.

\begin{lstlisting}[caption={Policy declaration in \texttt{aisafepy.flow}.}, label={lst:policy}]
policy = (
    Policy()
    .label_source("user_prompt", integrity="TRUSTED")
    .label_source("gmail.read",  integrity="UNTRUSTED",
                  caps={Capability.READ_USER})
    .label_source("web.fetch",   integrity="UNTRUSTED")
    .require("send_email",     control_flow_integrity="TRUSTED",
             caps={Capability.WRITE_EXTERNAL})
    .require("payments.transfer", control_flow_integrity="TRUSTED",
             caps={Capability.WRITE_EXTERNAL})
    .deny_if("send_email",
             when=lambda to, body: "read.secrets" in body.provenance,
             reason="secret-to-external-sink")
    .declassify("summarize", by="P-LLM only, structured output")
)
\end{lstlisting}

\paragraph{Enforcement modes.} Two modes. In \texttt{strict} mode,
denial raises \texttt{IFCViolationError}, following CaMeL's
deny-by-construction model. In \texttt{mediated} mode, denial routes
to a human-approval callback before deciding, trading a small
fraction of utility for very high recall.

\paragraph{Adapters.} A single shim per framework wires the policy
into the host runtime. For OpenAI Agents SDK we monkey-patch each
\texttt{FunctionTool}'s invocation callable with a check-then-call
proxy. For LangGraph we wrap nodes at registration time and ferry
the IFC context through the state dict. For LlamaIndex we replace
the underlying \texttt{fn} on each tool. For Anthropic we adapt the
user-supplied dispatch handler. The MCP adapter is the most
elaborate: it patches \texttt{add\_tool} to perform static
description screening (regex over a community-curated pattern set
of known MCP tool-poisoning markers) and hashed-manifest comparison,
then patches \texttt{handle\_call\_tool} for runtime IFC.

\subsection{aisafepy.stream: cascaded streaming guardrails}

The pipeline runs Tier 1 (deterministic, $<$10 ms), Tier 2 (small
classifier or forecaster, $<$100 ms), and Tier 3 (heavyweight LLM
judge or activation probe, $<$500 ms) in sequence. Tier 2 advances
by default after Tier 1. Tier 3 runs only when Tier 2's confidence
falls inside a configurable uncertainty band, typically
$[0.30, 0.70]$. Each tier runs all of its guards in parallel and
aggregates the results: any block wins, then any escalation, then
any transform, then the most confident allow.

\paragraph{Latency budget.} A \texttt{budget\_ms\_p95} parameter
sets a deadline. When the deadline is exceeded mid-cascade, the
pipeline emits a budget-exceeded \texttt{GuardDecision} with an
\texttt{evidence.budget\_exceeded} flag and either fails open
(default) or fails closed, configurable via \texttt{on\_budget\_exceeded}.
This makes the safety/UX tradeoff explicit instead of hiding it.

\paragraph{Activation probes.} For self-hosted models, the
\texttt{aisafepy.stream.probes} subpackage exposes
\texttt{LinearProbe} and \texttt{MLPProbe} classes that fit
logistic-regression or two-layer-MLP probes over selected hidden
states. The fit procedure follows Anthropic's recipe
\citep{anthropicprobes}: mean-pool activations across token
positions at a chosen set of layers, concatenate, fit. The probe
header is about 50 KB on disk and runs in under 2 ms on commodity
GPUs.

\paragraph{DoS defenses.} Every Tier-1 regex guard has a
\texttt{max\_input\_chars} cap (default 50,000) to defend against
adversarial false-positive vectors that drive content classifiers
to block large fractions of benign traffic via short, crafted
inputs. The same cap protects against $O(n^2)$ regex
catastrophic-backtracking attacks.

\subsection{aisafepy.adapt: continuous eval-to-guardrail compilation}

The compiler ingests a stream of \texttt{FailureRecord}s and
produces a \texttt{CompilationReport} of typed artifacts. The
pipeline is:

\begin{enumerate}[leftmargin=2em, itemsep=0.25em]
  \item \textbf{Sources.} Adapters for PyRIT (\texttt{PyRITSource}),
    Garak (\texttt{GarakReport}), Inspect (\texttt{InspectLog}),
    and production traces (\texttt{LangfuseTraces}) normalize their
    native record formats into \texttt{FailureRecord}.

  \item \textbf{Clustering.} Encode failures with
    sentence-transformers, reduce with UMAP, cluster with HDBSCAN.
    Fallbacks to k-means or TF-IDF + agglomerative clustering when
    heavier dependencies are absent.

  \item \textbf{Per-cluster compilation.} For each cluster above
    the \texttt{min\_attack\_success\_rate} threshold, emit a
    \texttt{CompiledArtifact} for each declared target. Five
    targets are supported:
    \begin{itemize}[leftmargin=1.5em]
      \item \texttt{distill\_classifier}: fine-tune a small encoder
        (default Llama-Prompt-Guard-2-22M) on the cluster's
        positives plus contrastive negatives.
      \item \texttt{synthesize\_regex}: longest-common-substring
        search plus precision estimation against negatives,
        retaining the top-$k$ patterns above a precision threshold.
      \item \texttt{policy\_rule}: emit Cedar or OPA Rego rules for
        governance integrations.
      \item \texttt{steering\_vector}: contrastive activation
        steering (mean-difference or conditional, per ICLR 2025)
        producing a \texttt{.safetensors} file per cluster.
      \item \texttt{deliberative\_case}: append the cluster's
        prototypical examples to a structured policy prompt for the
        agent to consult before answering.
    \end{itemize}

  \item \textbf{Canary and governance.} \texttt{promote()} attaches
    the new guards to a \texttt{GuardPipeline} in shadow or active
    mode, gated by a false-positive budget. An append-only audit
    log with SHA-256 hash chaining records every compile and
    promote event.
\end{enumerate}
